\documentclass{intsys}
\usepackage{listings}
\usepackage{xcolor}
\usepackage[ruled,vlined]{algorithm2e}

\definecolor{codegreen}{rgb}{0,0.6,0}
\definecolor{codegray}{rgb}{0.5,0.5,0.5}
\definecolor{codepurple}{rgb}{0.58,0,0.82}
\definecolor{backcolour}{rgb}{0.95,0.95,0.92}

\lstdefinestyle{mystyle}{
  backgroundcolor=\color{backcolour},
  commentstyle=\color{codegreen},
  keywordstyle=\color{magenta},
  numberstyle=\tiny\color{codegray},
  stringstyle=\color{codepurple},
  basicstyle=\ttfamily\footnotesize,
  breakatwhitespace=false,
  breaklines=true,
  captionpos=b,
  keepspaces=true,
  numbers=left,
  numbersep=5pt,
  showspaces=false,
  showstringspaces=false,
  showtabs=false,
  tabsize=2,
  frame=single,
  literate=
  {а}{{\selectfont\char224}}1
  {б}{{\selectfont\char225}}1
  {в}{{\selectfont\char226}}1
  {г}{{\selectfont\char227}}1
  {д}{{\selectfont\char228}}1
  {е}{{\selectfont\char229}}1
  {ё}{{\"e}}1
  {ж}{{\selectfont\char230}}1
  {з}{{\selectfont\char231}}1
  {и}{{\selectfont\char232}}1
  {й}{{\selectfont\char233}}1
  {к}{{\selectfont\char234}}1
  {л}{{\selectfont\char235}}1
  {м}{{\selectfont\char236}}1
  {н}{{\selectfont\char237}}1
  {о}{{\selectfont\char238}}1
  {п}{{\selectfont\char239}}1
  {р}{{\selectfont\char240}}1
  {с}{{\selectfont\char241}}1
  {т}{{\selectfont\char242}}1
  {у}{{\selectfont\char243}}1
  {ф}{{\selectfont\char244}}1
  {х}{{\selectfont\char245}}1
  {ц}{{\selectfont\char246}}1
  {ч}{{\selectfont\char247}}1
  {ш}{{\selectfont\char248}}1
  {щ}{{\selectfont\char249}}1
  {ъ}{{\selectfont\char250}}1
  {ы}{{\selectfont\char251}}1
  {ь}{{\selectfont\char252}}1
  {э}{{\selectfont\char253}}1
  {ю}{{\selectfont\char254}}1
  {я}{{\selectfont\char255}}1
  {А}{{\selectfont\char192}}1
  {Б}{{\selectfont\char193}}1
  {В}{{\selectfont\char194}}1
  {Г}{{\selectfont\char195}}1
  {Д}{{\selectfont\char196}}1
  {Е}{{\selectfont\char197}}1
  {Ё}{{\"E}}1
  {Ж}{{\selectfont\char198}}1
  {З}{{\selectfont\char199}}1
  {И}{{\selectfont\char200}}1
  {Й}{{\selectfont\char201}}1
  {К}{{\selectfont\char202}}1
  {Л}{{\selectfont\char203}}1
  {М}{{\selectfont\char204}}1
  {Н}{{\selectfont\char205}}1
  {О}{{\selectfont\char206}}1
  {П}{{\selectfont\char207}}1
  {Р}{{\selectfont\char208}}1
  {С}{{\selectfont\char209}}1
  {Т}{{\selectfont\char210}}1
  {У}{{\selectfont\char211}}1
  {Ф}{{\selectfont\char212}}1
  {Х}{{\selectfont\char213}}1
  {Ц}{{\selectfont\char214}}1
  {Ч}{{\selectfont\char215}}1
  {Ш}{{\selectfont\char216}}1
  {Щ}{{\selectfont\char217}}1
  {Ъ}{{\selectfont\char218}}1
  {Ы}{{\selectfont\char219}}1
  {Ь}{{\selectfont\char220}}1
  {Э}{{\selectfont\char221}}1
  {Ю}{{\selectfont\char222}}1
  {Я}{{\selectfont\char223}}1
}
\lstset{style=mystyle}

\ShortTitle{Гибридное решение параметрических задач}
\TitleRU{
  Гибридное решение параметрических задач: координация логического
  вывода и внешних алгоритмов
}
\TitleEN{
  Hybrid Solution of Parametric Problems: Coordination of Logical
  Inference and External Algorithms
}
\Author{
  Анненков Александр Петрович % ФИО на русском
}{
  аспирант каф. математической теории интеллектуальных систем мех.-мат. ф-та МГУ
}{
  Annenkov Alexander Petrovich % ФИО на английском
}{
  graduate student, Lomonosov Moscow State University, Faculty of
  Mechanics and Mathematics, Chair
  of Mathematical Theory of Intellectual Systems
}{
  just.std@yandex.ru % e-mail
}{%
  0009-0002-9332-1569% ORCID
}
\AbstractRU{
  TODO
}
\AbstractEN{
  TODO
}
\KeywordsRU{
  автоматический вывод, делегирование, разбор случаев,
  параметрические задачи, системы компьютерной алгебры%
}
\KeywordsEN{
  automated reasoning, delegation, case analysis,
  parametric problems, computer algebra systems%
}
\begin{document}

\section{Архитектура механизмов делегирования}

В описываемой системе логические правила могут делегировать
вычисления двумя способами. Первый --- \emph{вызов подзадачи}: то
же ядро системы~\cite{IS-Annenkov2026} рекурсивно решает вспомогательную
задачу. Второй
--- \emph{вызов внешнего алгоритма}: управление передаётся
детерминированной функции, реализованной за пределами правил. Оба
механизма единообразно встраиваются в правило, которое остаётся
декларативной парой <<шаблон $\Rightarrow$ вывод>> с
ограничениями-термами~\cite{IS-Annenkov2026}, без условных конструкций,
циклов и явных передач управления; делегирование выражается тем,
что некоторые ограничения или подвыражения вывода содержат вызов
подзадачи либо внешнего алгоритма. Ниже описаны обе формы и приведён сквозной
пример.

\subsection{Вызов подзадачи из правила}

Базовый механизм подзадач описан в~\cite{IS-Annenkov2026}:
\emph{подзадача} --- вспомогательная задача, которая решается
отдельной сессией системы в рамках текущей задачи и допускает те же
типы целевых установок (\texttt{find}, \texttt{prove},
\texttt{transform}); её результат возвращается в основную задачу.
Подзадачи кэшируются, чтобы избежать повторного решения одинаковых
подзадач, и могут рекурсивно создавать собственные подзадачи с
ограничением глубины вложенности.

Настоящая работа дополняет этот механизм возможностью запустить из
ограничения правила подзадачу типа \texttt{find} на произвольно
заданном наборе переменных и условий. Ответ подзадачи связывается с
параметром правила и далее участвует в его выводе наравне с другими
связанными параметрами. Тем самым правило получает способ вычислить
вспомогательный результат, не выражая логику этого вычисления
самостоятельно --- её предоставляют другие правила, срабатывающие
уже внутри подзадачи.

При запуске такой подзадачи в её контекст переносятся известные
факты родительской задачи. Целевой терм родителя в подзадачу не
передаётся: у подзадачи своя цель --- поиск заданных переменных при
заданных условиях, и решается она независимо от пути решения
основной задачи.

Защита от бесконечной рекурсии встроена в сам механизм кэша
подзадач~\cite{IS-Annenkov2026}: подзадача записывается в кэш до
начала своего решения с пометкой <<решается>>. Если в ходе её
решения возникает запрос на ту же подзадачу, кэш вместо запуска
нового сеанса возвращает маркер <<решение не найдено>>, и
рекурсивная ветвь обрывается.

Рекурсивный вызов возникает, например, когда правило применимо к
собственному результату: одна из его ветвей делегирует подзадачу,
в условиях которой снова срабатывает то же самое правило. Если
структура подзадачи в точности совпадает со структурой исходной
задачи, кэш обрывает её на первом же повторении. Если же
делегируемая подзадача строго проще исходной --- например, правило
понижает степень многочлена и делегирует уравнение меньшей степени,
либо фиксирует значение одного из параметров и делегирует задачу с
уже подставленным значением, --- то структура подзадачи отличается,
она получает свежую запись в кэше и решается нормально. Каждое
последующее вхождение работает с более простой задачей, поэтому
рекурсия в таких случаях конечна сама по себе и защита кэша до неё
не доходит. Пример такого редуцирующего вхождения --- карточка C3
в~\S2.3.

Конкретный синтаксис записи подзадачи в требованиях правила и пример
сочетания такого вызова с обращением к внешнему алгоритму приведены
в~\S2.3 на карточке C3.

\subsection{Вызов внешнего алгоритма из правила}

Параллельно механизму подзадач описываемая система поддерживает
обращение к \emph{внешнему алгоритму} --- детерминированной функции,
реализованной за пределами логического слоя; например, встроенной в
саму систему или обёртывающей вызов внешней системы компьютерной
алгебры. Такой алгоритм оборачивается в \emph{процедурный
символ} --- символ произвольной арности, результат применения которого
к аргументам вычисляется этим алгоритмом в процессе нормализации.
Функцию, реализующую процедурный символ, далее будем называть
\emph{калькулятором}. Тем же механизмом реализованы и базовые
арифметические операции: вычисление константных подвыражений входит
в процедуру нормализации терма~\cite{IS-Annenkov2026} и выполняется
встроенным калькулятором.

Правило не различает процедурный и обычный символ: процедурный
символ может появляться в шаблоне, в ограничениях правила или в его
выводе наравне с обычным. Никаких особых конструкций для вызова
калькулятора в языке правил не предусмотрено.

Процедурный символ в шаблоне правила не вычисляется --- сопоставление
с шаблоном работает структурно, как с любым другим символом.
Вычисление калькулятора происходит позже, при нормализации
ограничения или вывода после подстановки параметров правила.

Калькулятор принимает термы --- аргументы своего процедурного
символа --- и возвращает терм. Допускается \emph{отказ от вычисления}:
калькулятор может вернуть исходный терм без изменений, оставив
процедурный символ в выражении. В остальных случаях возвращённый терм
занимает место вызова в окружающем выражении ограничения правила;
именно структура этого выражения и определяет, как правило использует
результат вызова.

Терм с вызовом калькулятора может связать параметр правила с
конкретным значением, перечислить множество его допустимых значений
или выступить в роли предиката. Это определяется тем, свободен или
связан параметр правила в позиции, унифицируемой с результатом
вызова. \emph{Свободным} называется параметр правила, для которого
в текущей попытке применения правила ещё не подобрана подстановка;
\emph{связанным} --- параметр, для которого подстановка уже
зафиксирована, например из шаблона правила или ранее обработанного
ограничения.

Если параметр свободен, возвращённый терм поставляется ему как
подстановка, а множество с несколькими элементами раскрывает ветви
декартовым произведением. Это естественное обобщение стандартного
механизма поиска подстановок параметров из~\cite{IS-Annenkov2026}:
ограничения вида $p = t$ итеративно просматриваются, и новые
связанные параметры раскрывают дополнительные ограничения на
следующей итерации. Если же параметр в позиции, унифицируемой с
результатом, уже связан или там стоит константа, ограничение
нормализуется и проверяется на истинность; ветви, на которых оно не
выполняется, отбрасываются.

Конкретные синтаксические формы обоих режимов показаны в сквозном
примере~\S1.3.

\subsection{Пример}

В качестве сквозного примера рассмотрим правило поиска рациональных
корней многочлена с целыми коэффициентами. Теоретическое основание ---
классическая теорема: монический многочлен с целыми коэффициентами
имеет рациональным корнем только целочисленный делитель свободного
члена. Правило формализует следствие теоремы как редукционный шаг:

\begin{lstlisting}
rule {
    attr level(5), goal(find(x, ..));

    F == 0 =>
        (x == u) || (Q == 0);

    is_polynomial(F, x),
    d == free_term(F, x),
    d != 0,
    u in divisors(d),
    substitute(x == u, F) == 0,
    Q == poly_div(F, x - u);
}
\end{lstlisting}

Шаблон $F = 0$ связывает параметр правила $F$ с многочленом из
условий задачи. Вывод правила $(x = u) \mathrel{||} (Q = 0)$ ---
дизъюнкция двух альтернатив: левая ветвь фиксирует найденный
корень, правая редуцирует задачу к многочлену меньшей степени. В
ограничениях задействованы пять
процедурных символов: \texttt{is\_polynomial}, \texttt{free\_term},
\texttt{divisors}, \texttt{substitute}, \texttt{poly\_div}. Их
реализация неоднородна: часть встроена в ядро самой системы, часть
--- обёртки над SymPy, использующие CAS как вспомогательный механизм
для отдельных полиномиальных операций. С точки зрения правила это
различие не существенно: все пять --- процедурные символы с единым
интерфейсом, и правило обрабатывает их единообразно.

Применение правила к задаче $F = x^3 - 6x^2 + 11x - 6$, $F = 0$, где
$x$ --- искомое:

\begin{enumerate}
  \item Шаблон $F = 0$ связывает $F$ с $x^3 - 6x^2 + 11x - 6$.
  \item \texttt{free\_term(F, x)} вычисляет свободный член,
    связывая $d \mapsto -6$. Требование $d \neq 0$ выполнено.
  \item \texttt{divisors(d)} выдаёт множество
    $\{\pm 1, \pm 2, \pm 3, \pm 6\}$. Связывание $u \in
    \texttt{divisors}(d)$ раскрывает восемь ветвей.
  \item В каждой ветви \texttt{substitute(x == u, F)} вычисляется и
    сравнивается с нулём; так как $u$ уже связан, ограничение
    работает как проверка истинности. Ветви $u = 1, 2, 3$ проходят;
    остальные отсеиваются.
  \item Для каждой прошедшей ветви \texttt{poly\_div(F, x - u)}
    вычисляет частное $Q$. Дизъюнктивная резолюция
    $(x = u) \mathrel{||} (Q = 0)$ формирует кусочный ответ:
    три ветви с зафиксированным корнем плюс редуцированная задача
    для дальнейшей обработки.
\end{enumerate}

Подзадачи в смысле~\S1.1 здесь не используются явно: одно правило
целиком выполняет один шаг редукции. В~\S2 разобран случай, когда
правило в одной из ветвей делегирует решение полностью внешнему
алгоритму \texttt{sympy\_solve}.

\section{Гибридный подход при решении задач}

Механизмы делегирования из~\S1 нейтральны: они применимы к любому
классу задач. В данной работе мы рассмотрим конкретный класс, в
котором делегирование даёт качественный выигрыш, ---
параметрические задачи, на которых системы компьютерной алгебры
систематически теряют решения. Настоящий раздел показывает, как
описанные механизмы работают в этом классе.

Раздел организован в четыре части. В \S2.1 зафиксированы
проблематика и каталог репрезентативных задач. В \S2.2 описано
семейство логических приёмов, складывающихся в гибридную схему. В
\S2.3 приведены представительные примеры. В \S2.4 подведены итоги
по представительным примерам.

\subsection{Классы параметрических задач, вызывающих трудности у SymPy}

Современные системы компьютерной алгебры, на примере SymPy~1.14.0
\cite{sympy}, систематически теряют решения в задачах с явным
параметром. Например, \texttt{SymPy.solve(Eq(a*x, a), x)} возвращает
$[1]$, теряя случай $a = 0$, в котором уравнение становится
тождеством. Типичный механизм потери --- деление на содержащее
параметр выражение без явного отслеживания случая его обнуления.
Аналогичные эффекты наблюдаются и в других местах: при работе с
радикалами теряются ОДЗ-условия; для модулей не реализован разбор
по знаку аргумента.

В качестве первичного материала для исследования собран небольшой
каталог верифицированных параметрических задач; каждая
сопровождается минимальным SymPy-скриптом, демонстрирующим потерю
или искажение решения. Каталог используется ниже как источник
представительных примеров (\S2.3) и опорная точка для систематизации
логических приёмов (\S2.2).

\subsection{Решение задач с параметрами с использованием гибридного подхода}

Гибридная схема состоит из логических преобразований и обращений к
внешнему алгоритму. Логический слой представлен семейством приёмов,
которые в одном правиле могут свободно сочетаться. Типичные приёмы:

\begin{itemize}
  \item \emph{Факторизация выражения} --- переписывание в
    произведение, открывающее перспективу разбора случаев по
    обнулению множителей.
  \item \emph{Разбор случаев} --- переписывание задачи в виде
    конечного набора ветвей, каждая из которых далее обрабатывается
    отдельно. В коде правила это выражается дизъюнктивной формой
    вывода: вывод записывается как дизъюнкция нескольких
    альтернатив, и каждая ветвь даёт свою часть кусочного ответа.
    Сами ветви порождаются либо связыванием параметра правила
    (\S1.2) с неединичным множеством, либо явным перечислением
    альтернатив, когда их число известно заранее. Типичный частный
    случай --- раскрытие модулей: модуль $|t|$ порождает две ветви
    по знаку $t$, в которых $|t|$ заменяется на $t$ или $-t$.
  \item \emph{ОДЗ-фильтр кандидатов} --- ОДЗ-предикат проверяется на
    каждой ветви ответа внешнего алгоритма; ветви, не удовлетворяющие
    ОДЗ, отбрасываются.
  \item \emph{Редукция вопроса о разрешимости в зависимости от
    параметра} --- задача вида <<при каких значениях параметра
    исходная задача разрешима>> переписывается в обычную задачу
    поиска допустимых значений параметра --- например, через явное
    решение исходной задачи или анализ её экстремума, --- и далее
    решается обычным делегированием.
\end{itemize}

\subsection{Примеры}

\paragraph{C1 --- параметрическое линейное уравнение
$a \cdot x = a$.} \texttt{SymPy.solve(Eq(a*x, a), x)} возвращает
$[1]$, теряя случай $a = 0$, в котором уравнение становится
тождеством. Гибридный путь --- разбор случаев по обнулению
коэффициента с делегированием SymPy в ветви $a \neq 0$.
В реализации правило выглядит так:

\begin{lstlisting}
rule {
    attr level(2), goal(find(x, ..));

    x*a + b == 0 =>
        answer(
            (a == 0 && b == 0 && x in Reals) ||
            (a == 0 && b != 0 && x in empty_set) ||
            (a != 0 && x in sympy_solve(x*a + b == 0, x))
        );

    a is known,
    b is known;
}
\end{lstlisting}

Из режимов взаимодействия~\S1.2 здесь представлены оба: связывание
параметров правила $a$ и $b$ из шаблона задаёт подстановку, а
предикатные калькуляторы \texttt{a is known} и \texttt{b is known}
проверяют истинность ограничения уже на связанных значениях. Случаи
перечислены явно в выводе правила, записанном как дизъюнкция внутри
\texttt{answer}; связывание параметра по неединичному множеству
здесь не используется --- число ветвей известно заранее.
Делегирование в ветви $a \neq 0$
обращается к процедурному символу \texttt{sympy\_solve}: ветвь
активируется только после того, как доказано $a \neq 0$, и SymPy
возвращает корректное решение $x = 1$.

Итог: общий ответ собирается как
$(a = 0 \,\&\&\, x \in \mathbb{R}) \mathrel{||} (a \neq 0 \,\&\&\,
x = 1)$.

\paragraph{C2 --- параметрическое линейное уравнение
$(a^2 - 4) \cdot x = a - 2$.} \texttt{SymPy.solve(\ldots, x)} возвращает
$1/(a + 2)$, теряя случай $a = 2$, где обе части обнуляются и
решение --- всё $\mathbb{R}$, и случай $a = -2$, где левая часть
равна нулю, а правая --- нет, и решений не существует. Здесь
срабатывает то же универсальное правило линейного уравнения из
листинга C1: после факторизации $a^2 - 4 = (a-2)(a+2)$ коэффициенты
становятся известными, и разбор случаев по обнулению старшего
коэффициента даёт правильный кусочный ответ. Карточка показывает,
что одно правило переиспользуется для целого семейства
параметрических задач с разной формой коэффициентов.

\paragraph{C3 --- параметрическое квадратное уравнение
$(a^2 - 1) \cdot x^2 + 2(a - 1) \cdot x + 1 = 0$.} \texttt{solve(\ldots)}
возвращает пару корней с $(a^2 - 1)$ в знаменателе и $2 - 2a$ под
корнем; при $a = 1$ обе формулы вырождаются в $0/0$, при $a = -1$
знаменатель обнуляется. SymPy не разбирает старший коэффициент и не
замечает, что при $a = \pm 1$ тип уравнения меняется --- причём при
$a = 1$ одновременно обнуляется и линейный коэффициент. Гибридный
путь --- двухуровневый разбор случаев на универсальном правиле
квадратного уравнения $a \cdot x^2 + b \cdot x + c = 0$:

\begin{lstlisting}
rule {
    attr level(2), goal(find(x, ..));

    a*x^2 + b*x + c == 0 =>
        answer(
            (a == 0 && L) ||
            (a != 0 && x in sympy_solve(a*x^2 + b*x + c, x))
        );

    a is known,
    b is known,
    c is known,
    L == find(x) { b*x + c == 0; b is known; c is known };
}
\end{lstlisting}

В ветви $a \neq 0$ правило делегирует уравнение целиком символу
\texttt{sympy\_solve} в духе \S1.2. В ветви $a = 0$ ограничение
\texttt{L == find(x) \{ b*x + c == 0; \ldots \}} запускает
подзадачу \S1.1 на линейном уравнении $b \cdot x + c = 0$ с
известными $b$ и $c$; внутри неё срабатывает то же правило
линейного разбора случаев, что в листинге C1, и его ответ
связывается с параметром \texttt{L}. Применение правила к C3 даёт
кусочный ответ
\[
  (a = 1 \wedge x \in \emptyset)
  \mathrel{||}
  (a = -1 \wedge x = 1/4)
  \mathrel{||}
  \bigl(a \neq \pm 1 \wedge x \in
  \tfrac{- a + 1 \pm \sqrt{2 - 2a}}{a^2 - 1}\bigr).
\]
Ветви $a = \pm 1$ возникают рекурсивно: при $a = 1$ подзадача
получает $b = 0, c = 1$ → противоречие → $\emptyset$; при $a = -1$ ---
$b = -4, c = 1$ → $x = 1/4$ через \texttt{sympy\_solve}. Карточка
демонстрирует одновременное использование обоих механизмов
делегирования из \S1: подзадача обрабатывает вырожденную ветвь
после понижения степени уравнения до линейной, а внешний CAS ---
невырожденную; логика разбора по обнулению старшего коэффициента
остаётся на стороне правила.

\subsection{Результаты}

В реализованной системе гибридный подход даёт корректный кусочный
ответ на каждой из трёх представительных карточек C1, C2, C3 ---
задачах, на которых SymPy~1.14.0 теряет часть решений. На карточке
C3 продемонстрировано одновременное использование обоих механизмов
делегирования из~\S1: подзадача обрабатывает вырожденную ветвь,
внешний CAS --- невырожденную, а разбор случаев по старшему
коэффициенту остаётся на стороне правила.

\appendix

\section{Алгоритм конкретизации параметров гипотезы}
\label{appendix:ground}

\emph{Гипотеза}, введённая в~\cite{IS-Annenkov2026}, --- пара
$(\mathrm{resolution}, \mathrm{requirements})$ из правой части
правила после подстановки параметров и списка оставшихся требований.
Для процедуры конкретизации удобно расширить это понятие до кортежа
$(\mathrm{resolution}, \mathrm{requirements}, \mathrm{params},
\mathrm{free\_params})$, дополнительно отделяя уже связанные параметры
правила (\texttt{params}) от ещё не связанных
(\texttt{free\_params}).

В данном приложении приводится псевдокод процедуры конкретизации
параметров гипотезы, реализующей механизм связывания параметров,
описанный в~\S1.2: она переводит гипотезу со свободными параметрами
в одну или несколько конкретизированных, последовательно применяя
два механизма --- связывание параметров через требования вида
$t == p$ и раскрытие декартова произведения по генераторам
$p \in \mathrm{set}(\ldots)$. Конкретизация заканчивается успехом на
тех ветвях, на которых после связывания параметров значениями из
генераторов и повторного применения $\mathrm{BindEqualityParams}$
множество \texttt{free\_params} становится пустым; ветви, в которых
остался хотя бы один свободный параметр, отбрасываются.

\paragraph{Вспомогательный метод $\mathrm{BindEqualityParams}$.} Просматривает
\texttt{requirements} в поисках требований вида $p == v$ (или
$v == p$), где $p \in \mathrm{free\_params}$. Из всех таких
требований выбирается то, у которого $v$ не содержит свободных
параметров; выполняется подстановка $p \mapsto v$ во всю гипотезу,
после чего соответствующее требование удаляется. Шаг повторяется,
пока такие требования есть. Оставшиеся $p == v$ с параметрическим
$v$ возвращаются в \texttt{requirements} как требования к будущей
резолюции. По окончании производится повторная нормализация
требований и отбрасывание тривиально истинных.

\paragraph{Вспомогательный метод $\mathrm{ExtractGenerators}$.} Извлекает из
\texttt{requirements} требования вида
$p \in \mathrm{set}(e_1, \ldots, e_k)$, где
$p \in \mathrm{free\_params}$, и возвращает список пар
$(p, [e_1, \ldots, e_k])$. Соответствующие требования удаляются из
гипотезы; элементы $e_i$ становятся источником ветвей при последующем
раскрытии декартова произведения.

\begin{algorithm}[H]
  \caption{Конкретизация параметров гипотезы}
  \label{alg:ground}
  \footnotesize
  \SetKwInOut{Input}{Вход}
  \SetKwInOut{Output}{Выход}

  \Input{Гипотеза $H = (\mathrm{resolution}, \mathrm{requirements},
  \mathrm{params}, \mathrm{free\_params})$}
  \Output{Список конкретизированных гипотез $\mathcal{G}$}

  $\mathrm{BindEqualityParams}(H)$ \;
  \If{$H.\mathrm{free\_params} = \varnothing$}{
    \Return{$[\,H\,]$} \;
  }

  $\Gamma \gets \mathrm{ExtractGenerators}(H)$ \;
  \If{$\Gamma = \varnothing$}{
    \Return{$[\,]$} \;
  }

  $\mathcal{G} \gets [\,]$ \;
  \ForEach{$\sigma \in \prod\limits_{(p,\, S) \in \Gamma}
  \{\, p \mapsto e \mid e \in S \,\}$}{
    $H' \gets \mathrm{clone}(H)$ \;
    $\mathrm{substitute}(H', \sigma)$ \;
    $\mathrm{BindEqualityParams}(H')$ \;
    \If{$H'.\mathrm{free\_params} = \varnothing$}{
      $\mathcal{G} \gets \mathcal{G} \cup \{H'\}$ \;
    }
  }

  \Return{$\mathcal{G}$} \;
\end{algorithm}
\vfill

\begin{thebibliographyRU}{9}
  \RBibitem{IS-Annenkov2026}
  \by А.\,П.~Анненков
  \paper Архитектура упрощённой системы автоматического вывода на
  основе подхода А.\,С.~Подколзина
  \jour Интеллектуальные системы. Теория и приложения
  \yr 2026
  \vol 30
  \issue 1
  \pages 6--35

  \RBibitem{Podkolzin2008}
  \by А.\,С.~Подколзин
  \book Компьютерное моделирование логических процессов
  \bookvol 1
  \voltitle Архитектура и языки решателя задач
  \publaddr Москва
  \publ Физматлит
  \yr 2008

  \RBibitem{sympy}
  \by A.~Meurer, C.\,P.~Smith, M.~Paprocki et al.
  \paper SymPy: symbolic computing in Python
  \jour PeerJ Computer Science
  \yr 2017
  \vol 3
  \pages e103
  \crossref{https://doi.org/10.7717/peerj-cs.103}

  \RBibitem{Malpas1987}
  \by Дж.~Малпас
  \book Реляционный язык пролог и его применение
  \publaddr Москва
  \publ Наука
  \yr 1990

  \RBibitem{Mendelson1964}
  \by Э.~Мендельсон
  \book Введение в математическую логику
  \publaddr Москва
  \publ Наука
  \yr 1971

  \RBibitem{BaaderNipkow1998}
  \by F.~Baader, T.~Nipkow
  \book Term Rewriting and All That
  \publaddr Cambridge
  \publ Cambridge University Press
  \yr 1998
  \crossref{https://doi.org/10.1017/CBO9781139172752}

\end{thebibliographyRU}

\begin{thebibliographyEN}{9}
  \Bibitem{IS-Annenkov2026}
  \by A.\,P.~Annenkov
  \paper Architecture of a Simplified Automated Reasoning System
  Based on A. S. Podkolzin's Approach
  \jour Intelligent Systems. Theory and Applications
  \yr 2026
  \vol 30
  \issue 1
  \pages 6--35
  \lang In Russian

  \Bibitem{Podkolzin2008}
  \by A.\,S.~Podkolzin
  \book Computer modeling of logical processes
  \bookvol 1
  \voltitle Architecture and languages of the problem solver
  \publaddr Moscow
  \publ Fizmatlit
  \yr 2008

  \Bibitem{sympy}
  \by A.~Meurer, C.\,P.~Smith, M.~Paprocki et al.
  \paper SymPy: symbolic computing in Python
  \jour PeerJ Computer Science
  \yr 2017
  \vol 3
  \pages e103
  \crossref{https://doi.org/10.7717/peerj-cs.103}

  \Bibitem{Malpas1987}
  \by J.~Malpas
  \book Prolog: A Relational Language and Its Applications
  \publaddr Englewood Cliffs, N.\,J.
  \publ Prentice-Hall
  \yr 1987

  \Bibitem{Mendelson1964}
  \by E.~Mendelson
  \book Introduction to Mathematical Logic
  \publaddr Princeton, N.\,J.
  \publ Van Nostrand
  \yr 1964

  \Bibitem{BaaderNipkow1998}
  \by F.~Baader, T.~Nipkow
  \book Term Rewriting and All That
  \publaddr Cambridge
  \publ Cambridge University Press
  \yr 1998
  \crossref{https://doi.org/10.1017/CBO9781139172752}
\end{thebibliographyEN}

\end{document}
